\documentclass[12pt]{article}
\usepackage{graphicx} % Required for inserting images
\usepackage[margin=0.5in]{geometry}
\usepackage{amsmath, amssymb}
\usepackage{mathrsfs}


\usepackage{soul}


% Dr. David Brown Commands
\newcommand{\ds}{\displaystyle}
\newcommand{\R}{\mathbb{R}}
\newcommand{\N}{\mathbb{N}}
\newcommand{\Q}{\mathbb{Q}}
\newcommand{\C}{\mathbb{C}}


% Commands to get script r
\usepackage{calligra}
\DeclareMathAlphabet{\mathcalligra}{T1}{calligra}{m}{n}
\DeclareFontShape{T1}{calligra}{m}{n}{<->s*[2.2]callig15}{}
\newcommand{\scripty}[1]{\ensuremath{\mathcalligra{#1}}}

% Unit commands
\newcommand{\degree}{\text{°}}
\newcommand{\unitSpeed}{\frac{\text{m}}{\text{s}}}
\newcommand{\unitAcceleration}{\frac{\text{m}}{\text{s}^2}}

% Constant commands
\newcommand{\avogadro}{6.022\times10^{23}}

\newcommand{\boltzmannConstK}{1.381\times10^{-23}\frac{\text{J}}{\text{K}}}
\newcommand{\planckConst}{6.626\times10^{-34}\text{J}\cdot\text{s}}

\newcommand{\atmP}{1.01325\times10^5\,\text{Pa}}

% Commands to easily write partial derivatives
\newcommand{\partDeriv}[2]{\frac{\partial #1}{\partial #2}}
\newcommand{\doublePartDeriv}[2]{\frac{\partial^2 #1}{\partial^2 #2}}


\title{Problem 1.34}
\author{Gabriel Decker}


\begin{document}



\maketitle
\begin{flushleft}

Note that it seems like in my version of the textbook, processes C and D are swapped from what you'd expect. I'm going off of what the book actually says, not what you'd expect.\\~\\

For this problem, we have a diatomic gas.
At room temp, the vibrational degrees of freedom are frozen out, so use $f=5$ degrees of freedom (3 translational, 2 rotational).\\~\\

We'll need the following equations.
    First the equipartition theorem (used with $f=5$ as discussed):
\begin{equation}
    U_\text{thermal}=\frac{1}{2}NfkT
\end{equation}
then the equation of work done on quasistatic compression-work-relevant systems in terms of $P$ and $V$:
\begin{equation}
    W=-\int_{V_i}^{V_f}P(V)\,dV
\end{equation}
then the first law of thermodynamics:
\begin{equation}
    \Delta U=W+Q
\end{equation}
and finally the following version of the ideal gas law.
\begin{equation}
    PV=NkT
\end{equation}\\~\\

We can combine the equipartition thm and ideal gas law to get an equation for our thermal energy that is more useful to us.
$$U_\text{thermal}=\frac{5}{2}PV$$
When we change the volume or pressure, this will cause $U_\text{thermal}$ to change.
So, we use the following equations to calculate this change, depending on the situation.
$$\Delta U_\text{thermal}=\frac{5}{2}\Delta PV$$
$$\Delta U_\text{thermal}=\frac{5}{2}P\Delta V$$\\~\\

\textbf{Part a}\\
We'll run through all the steps A-D. We'll start by finding $W$.
Then we'll use the equipartition thm to find $\Delta U$.
Finally, we'll calculate $Q$.\\~\\

\textit{Process A:}\\
In step A, we are at a constant volume $V_1$ and are increasing the pressure from $P_1$ to $P_2$.
This means that $\Delta V=0$ and $\Delta P=P_2-P_1$.
Since $V_i=V_f$, this implies that $\ds W=-\int_{V_1}^{V_1}P\,dV=0$.
Thus, $\Delta U=Q$.
By our previously derived versions of $\Delta U$, this implies that
$$\Delta U_\text{thermal}=\frac{5}{2}\Delta PV$$
$$\implies\Delta U_\text{thermal}=\frac{5}{2}(P_2-P_1)V_1$$
Since $\Delta U=Q$,
$$\implies Q=\frac{5}{2}(P_2-P_1)V_1$$\\~\\

\textit{Process B:}\\
In step B, we are at a constant pressure $P_2$ and are increasing the volume from $V_1$ to $V_2$.
This means that $\Delta P=0$ and $\Delta V=V_2-V_1$.
Plugging in what we have into the equation for $W$, we get the following.
$$W=-\int_{V_i}^{V_f}P(V)\,dV$$
$$\implies W=-\int_{V_1}^{V_2}P_2\,dV$$
$$\implies W=-P_2(V_2-V_1)$$
By our previously derived versions of $\Delta U$, this implies that
$$\Delta U_\text{thermal}=\frac{5}{2}P\Delta V$$
$$\implies\Delta U_\text{thermal}=\frac{5}{2}P_2(V_2-V_1)$$
Since $\Delta U=W+Q$,
$$\implies Q=\Delta U-W$$
$$\implies Q=\frac{5}{2}P_2(V_2-V_1)+P_2(V_2-V_1)$$
$$\implies Q=\frac{7}{2}P_2(V_2-V_1)$$\\~\\

\textit{Process C:}\\
In step C, we are at a constant pressure $P_1$ and are decreasing the volume from $V_2$ to $V_1$.
This means that $\Delta P=0$ and $\Delta V=-(V_2-V_1)$.
Plugging in what we have into the equation for $W$, we get the following.
$$W=-\int_{V_i}^{V_f}P(V)\,dV$$
$$\implies W=-\int_{V_2}^{V_1}P_1\,dV$$
$$\implies W=-P_1(V_1-V_2)$$
$$\implies W=P_1(V_2-V_1)$$
By our previously derived versions of $\Delta U$, this implies that
$$\Delta U_\text{thermal}=\frac{5}{2}P\Delta V$$
$$\implies\Delta U_\text{thermal}=-\frac{5}{2}P_1(V_2-V_1)$$
Since $\Delta U=W+Q$,
$$\implies Q=\Delta U-W$$
$$\implies Q=-\frac{5}{2}P_1(V_2-V_1)-P_1(V_2-V_1)$$
$$\implies Q=-\frac{7}{2}P_1(V_2-V_1)$$\\~\\

\textit{Process D:}\\
In step D, we are at a constant volume $V_2$ and are decreasing the pressure from $P_2$ to $P_1$.
This means that $\Delta V=0$ and $\Delta P=-(P_2-P_1)$.
Since $V_i=V_f$, this implies that $\ds W=-\int_{V_2}^{V_2}P\,dV=0$.
Thus, $\Delta U=Q$.
By our previously derived versions of $\Delta U$, this implies that
$$\Delta U_\text{thermal}=\frac{5}{2}\Delta PV$$
$$\implies\Delta U_\text{thermal}=-\frac{5}{2}(P_2-P_1)V_2$$
Since $\Delta U=Q$,
$$\implies Q=-\frac{5}{2}(P_2-P_1)V_2$$\\~\\~\\


\textbf{Part b}\\
In process A, no work is done (no piston compression or expansion), but the gas inside is being heated.
In process B, work is being done to the piston to expand the volume. The gas is being heated.
In process C, work is being done by the piston to constrict the volume. The gas is being cooled.
In process D, no work is done (no piston compression or expansion), but the gas inside is being cooled.\\~\\~\\


\textbf{Part c}\\
The net work, heat, and change in energy is just the sum of the previously computed values.
$$W_\text{tot}=0-P_2(V_2-V_1)+P_1(V_2-V_1)+0=-(P_2-P_1)(V_2-V_1)$$
$$Q_\text{tot}=\frac{5}{2}(P_2-P_1)V_1+\frac{7}{2}P_2(V_2-V_1)-\frac{7}{2}P_1(V_2-V_1)-\frac{5}{2}(P_2-P_1)V_2$$
$$\implies Q_\text{tot}=\frac{7}{2}(P_2-P_1)(V_2-V_1)-\frac{5}{2}(P_2-P_1)(V_2-V_1)$$
$$\implies Q_\text{tot}=(P_2-P_1)(V_2-V_1)=-W_\text{tot}$$
$$\Delta U_\text{tot}=\frac{5}{2}(P_2-P_1)V_1+\frac{5}{2}P_2(V_2-V_1)-\frac{5}{2}P_1(V_2-V_1)-\frac{5}{2}(P_2-P_1)V_2$$
$$\implies\Delta U_\text{tot}=-\frac{5}{2}(P_2-P_1)(V_2-V_1)+\frac{5}{2}(P_2-P_1)(V_2-V_1)=0$$\\~\\

Throughout this cycle, the net work done on the gas is negative.
So, the gas is doing work on the piston.
The net heat applied to the system is positive, so you're heating the gas, on average.
The net change in energy is zero, which also makes it so the fact that $W=-Q$ in this case make sense.





\end{flushleft}

\end{document}
